%
% Komplexe Zahlen und Funktionen
%
\subsubsection{Komplexe Zahlen und Funktionen}

%
% konstruktion: Konstruktion
%
Mit komplexen Zahlen lassen sich geometrische Konstruktionen in
\index{Konstruktion, geometrisch}%
der Ebene besonders einfach beschreiben.
\emph{Timon Gnehm}
\index{Gnehm, Timon}%
\index{Timon Gnehm}%
und
\emph{Pascal Hirzel}
\index{Pascal Hirzel}%
\index{Hirzel, Pascal}%
stellen in Kapitel~\ref{chapter:konstruktion} den Satz von Napoleon
\index{Satz!von Napoleon}%
vor.

%
% moebius: Möbius-Transformationen
%
Die Bewegungen der ebenen euklidischen Geometrie können in der 
komplexen Ebene breiter gefasst werden.
Indem die komplexe Ebene zur riemannschen Zahlenkugel erweitert wird,
\index{Zahlenkugel, riemannsch}%
können gebrochen rationale Transformationen, auch Möbius-Transformationen
\index{Mobius-Transformation@Möbius-Transformation}%
genannt, als Transformationsgruppe verwendet werden.
Sie sind zwar nicht mehr längentreu, aber immer noch winkeltreu.
\emph{Tom Dawson}
\index{Tom Dawson}%
\index{Dawson, Tom}%
zeigt in Kapitel~\ref{chapter:moebius}, wie man man Möbius-Transformationen
finden, analysieren und damit rechnen kann.
Insbesondere zeigt sich, dass die Zusammensetzung von Möbius-Transformationen
durch die Matrixmultiplikation beschrieben werden kann.

%
% geradlinig: Das seltsame Problem der geradlinigen Bewegung
%
Die Industrialisierung hat die Konstruktion von Maschinen nötig
gemacht, die eine Bewegung geradlinig führen können.
\emph{Shaarujan Kamalanathan}
\index{Shaarujan Kamalanathan}%
\index{Kamalanathan, Shaarujan}%
konstruiert die im 18.~und 19.~Jahrhundert gefundenen Apparate
in Kapitel~\ref{chapter:geradlinig} nach und zeigt, wie sie mit den
Eigenschaften der komplexen Zahlen verstanden werden können.

%
% julia: Julia-Mengen und Mandelbrot-Menge
%
Die Theorie der dynamischen System befasst sich mit dem Verhalten
von Iterationsfolgen.
Holomorphe Funktionen sind in diesem Zusammenhang besonders interessant,
da sie zu winkeltreuen Abbildungen der komplexen Ebene gehören.
\emph{Nicola Dall'Acqua} zeigt in Kapitel~\ref{chapter:julia}, wie
\index{Nicola Dall'Acqua}%
\index{Dall'Acqua, Nicola}%
die Julia- und Fatou-Mengen Punkte der komplexen Ebene anhand des
\index{Julia-Menge}%
\index{Fatou-Menge}%
Konvergenzverhaltens der zugehörigen Iterationsfolge klassifizieren.
Für holomorphe Funktionen entstehen filigrane selbstähnlichen Strukturen,
die diesen Mengen auch eine besondere ästhetische Qualität geben.
Die Mandelbrot-Menge im Parameterraum für quadratische Abbildungen
\index{Mandelbrot-Menge}%
wird ebenfalls eingeführt.

%
% jordan:
%
Der erste, aus heutiger Sicht unvollständige Beweis des Fundamentalsatzes
der Algebra durch Gauss scheiterte an ``inutitiv klaren'' Aussagen über
die Topologie der komplexen Ebene.
\emph{Dominik Castelberg}
\index{Dominik Castelberg}%
\index{Castelberg, Dominik}%
führt in Kapitel~\ref{chapter:jordan} durch die überraschenden
Schwierigkeiten eines vollständigen Beweises des jordanschen Kurvensatzes.
Er zeigt damit, dass zwischen einem ``anschaulichen'' Resultat und einem
strengen Beweis eine grosse Lücke klaffen kann.
Letztlich ist unsere Intuition insbesondere in höheren Dimensionen
unzuverlässig.

%
% qa: Quaternionen
%
Komplexe Zahlen vereinfachen die Beschreibung von Drehungen und Streckungen
in zwei Dimensionen.
Ist so etwas auch für Drehungen in drei Dimensionen möglich?
\emph{Kai Erdin}
\index{Kai Erdin}%
\index{Erdin, Kai}%
zeigt in Kapitel~\ref{chapter:qa}, dass die vierdimensionale
Quaternionenalgebra genau dies ermöglicht.
Allerdings ist die Formel für Drehungen viel komplizierter, da sie
der Tatsache Rechnung tragen muss, dass Drehungen im dreidimensionalen
Raum nicht vertauschen.

